\documentclass[10pt]{article}

\usepackage[utf8]{inputenc}

\usepackage[T1]{fontenc}

\usepackage{amsmath}

\usepackage{amsfonts}

\usepackage{amssymb}

\usepackage[version=4]{mhchem}

\usepackage{stmaryrd}

\usepackage{bbold}

\usepackage{graphicx}

\usepackage[export]{adjustbox}

\graphicspath{ {./images/} }

\usepackage{CJKutf8}



\begin{document}

савдера. Срезанные узлы (см. Узлов кобордизл) - это узлы внешнего рода нуль. Сигнатура и единицы Минковского узла определяются его классом кобордизмов. Функция на групле кобордизмов одномерных узлов в $S^{3}$ со значениями в группе $\mathbb{Z}$, к-рая сопоставляет классу кобордизмов сигнатуру представляющего его узла, является гомоморфизмом, образ к-рого совпадает с подгруппой четных чисел. Число заузливания всякого узла не менее половины его сигнатуры.



У. и з. к. Ф. тесно связаны с двулистными разветвленными накрывающими $\Sigma^{4}$ шара $D^{4}$ с ветвлением над ориентируемыми двумерными поверхностями $N \subset D^{4}$ с $N \cap S^{3}=\partial N=l$. В частности, сигнатура и единицы Минковского зацепления $L=\left(S^{3}, l\right)$ равны соответственно сигнатуре и единицам Минковского многообразия $\Sigma^{4}$. Край $\partial \Sigma^{4}=\Sigma^{3}$, будучи двулистным накрытием сферы $S^{3}$, разветвленным над зацеплением $l$, является инвариантом зацеплений $L$. В случае узлов $H_{1}\left(\Sigma^{3} ; \mathbb{Z}\right)$ является конечной группой. Эта группа, а также форма коэффициентов зацепления



\[

\lambda: H_{1}\left(\Sigma^{3} ; \mathbb{Z}\right) \otimes H_{1}\left(\Sigma^{3} ; \mathbb{Z}\right) \longrightarrow \mathbb{Q} / \mathbb{Z}

\]



определяется квадратичной формой узла следующим образом. Группой с зацеплением или $V$-г $\mathrm{p}$ у п ш о й наз. пара $(G, \mu)$, состоящая из конечной абелевой группы $G$ и невырожденной билинейной симметричной формы $\mu$ : $: G \otimes G \rightarrow \mathbb{Q} / \mathbb{Z}$. Каждой невырожденной симметричной целочисленной $(n \times n)$-матрице $A$ сопоставляется $V$ группа $(G, \mu)$, определяемая так: группа $G$ порождается әлементами $g_{1}, \ldots, g_{n}$ со следующими определяющими соотношениями: $\sum_{j=1}^{n} a_{i j} g_{j}=0, i=1, \ldots, n$, где $A=\left\|a_{i j}\right\|$, а $\mu\left(g_{i}, g_{j}\right)$ равно $\bmod 1$ элементу матрицы $A^{-1}$, находящемуся на месте $(i, j)$. Оказывается, что $V$-группа, определяемая этим способом по матрице $M+M^{\prime}$ квадратитной формы узла, изоморфна $V$-группе $\left(H_{1}\left(\Sigma^{3} ; \mathbb{Z}\right), \lambda\right)$ многообразия $\Sigma^{3}$ (см. [4], [9]). Числовые инварианты $V$-групп находятся методом Бланчффилда - Фокса [5]. С их помощью в нек-рых случаях удается различать узлы, имеющие изоморфыные группы.



Инварианты зацепления в двулистном накрытии $S^{3}$, разветвленном над узлом, могут быть найдены непосредственно по проекции узла с помощью следующей конструкции, к-рая приводит к к в а д р а т и ч н о й фо р м е ди а г р м мы узла. Регулярная проекция узла делит плоскость на области, к-рые могут быть однозначно окрашены в черный и белый цвета таким образом, чтобы бесконечная область $G_{0}$ была окрашена в черный цвет, а всякие две области, примыкающие друг к другу по дуге, были бы окрашены в разные цвета. Пусть $G_{0}, G_{1}, \ldots, G_{n}$ - все черные области. Каждой двойной точке $x$ диаграммы узла сопоставляется следующим образом нек-рое число $\eta(x) \in\{-1,0,1\}$. Пусть точка $x$ является общей граничной точкой двух черных областей $G_{i}$ и $G_{k}$. Если $i=k$, то $\eta(x)=0$. Если же $i \neq k$, то $\eta(x)=1$ тогда и только тогда, когда вращение от перехода к проходу вокруг точки $x$ по черной области происходит по часовой стрелке; в противном случае $\eta(x)=-1$. Образуется следующая $(n \times n)$-матрица $A=$ $=\left\|a_{i j}\right\|$, где $a_{i i}-$ сумма всех чисел $\eta(x)$, отвечающих двойным точкам $x$, лежацим на границе области $G_{i}$, a $a_{i k} \mathrm{c} i \neq k$ есть взятая с обратным знаком сумма всех чисел $\eta(x)$, где $x$ пробегает все общие граничные точки $G_{i}$ и $G_{k}$. Форма $f=\sum_{i, j=1}^{n} a_{i j} x_{i} x_{j}$ наз. к в а д р а т и чн о й $A=\left\|a_{i j}\right\|$ определяется типом узла с точностью до следующего отношения связанности: две квадратные матрицы наз. с в я 3 а н в ы м и, если одна может быть переведена в другую конечной последовательностью следующих операций: $Q_{1}: A \rightarrow T^{\prime} A T$, где $T-$ целочисленная унимодулярная матрица, $Q_{2}: A \rightarrow\|\|_{0 \pm 1}^{A} \|, \quad$ и обратных к ним. Модуль определителя матрицы $A$ является инвариантом узла и наз. д е т е р м и н а нт о м узла. Для всякого узла он нечетен и равен $|\Delta(-1)|$, где $\Delta(t)$ - многочлен Александера (см. Александера инварианты). V-группа, определяемая указаниым выше способом по матрице квадратичной формы произвольной диаграммы, является инвариантом узла. Более того, эта $V$-групша изоморфна $V$-группе $\left(H_{1}\left(\Sigma^{3}, \mathbb{Z}\right), \lambda\right)$ двулистного накрытия сферы $S^{3}$, разветвленного над узлом $K$.



\includegraphics[max width=\textwidth, center]{2023_07_09_78e3c4c8d363b6afcd59g-1(2)}

[2] G o e r i t z L, "MMath. Z.", 1933, Bd 36, S. 647; [3]' S e i-

f e r t H.. "Abh. Math. Sem. Hamb. Univ.", $1936, \mathcal{N}_{1} 1$, S. $84-$ $101 ;[4]$ K n es e r M., P u p p e D., "Math. Z.», 1953, Bd 58, S 376-84; [5] B l a n c h f i e l d R, F o x R., "Ann.' Math.", 1951, v. 53, p. 556-64; [6] T r o t t e r H , "Ann. Math.», 1962, S. 76, p. 464-68, [7] M u r a s u g i K , "Trans. Amer. Math. Phil.' Soc.», 1969, V. 66, p. $251-64$, [9] В и р о О. Я., "Изв. AH CCCP. Сер. матем.", 1973 , т. 37, с. $1241-58$ M. III. Фарбер



УЗЛОВ КОБОРДИЗМ (правильнее бордизм узлов см. Бор дизм) - отношение әквивалентности на множестве узлов, более слабое, чем изотопич. тип. Два гладких $n$-мерных узла $K_{1}=\left(S^{n+2}, k_{1}^{n}\right)$ и $K_{2}=\left(S^{n+2}, k_{2}^{n}\right)$ наз. к о б о р д а т т н м и, если существует гладкое ориентированное $(n+1)$-мерное подмногообразие $V$ многообразия $[0,1] \times S^{n+2}$, причем $V$ гомеоморфно $[0,1] \times S^{n}$ и $\partial V=V \cap[0,1] \times\left(S^{n+2}\right)=\left(0 \times k_{1}\right) \cup\left(1 \times-k_{2}\right)$. Здесь знак минус означает обращение ориентации. Узлы, кобордантные тривиальному узлу, наз. к о 6 о р-



\includegraphics[max width=\textwidth, center]{2023_07_09_78e3c4c8d363b6afcd59g-1(3)}

м и. Множество классов эквивалентности (кобордантности) $n$-мерных гладких узлов обозначается $C_{n}$. Oперация несвязной суммы определяет во множестве $C_{n}$ структуру абелевой группы. Обратным для класса У. к. $\left(S^{n+2}, k^{n}\right)$ служит класс У. к. $\left(-S^{k+2},-k^{n}\right)$. При всех четных $n$ группа $C_{n}$ равна нулю. Класс У.к. нечетномерного узла определяется его Зейферma матричей. Квадратная целочисленная матрица $A$ называется к о б о р д а т т о й н у л ю, если она унимодулярно конгруәнтна матрице вида $\left\|\begin{array}{ll}0 & N_{1} \\ N_{2} & N_{3}\end{array}\right\|$, где $N_{1}, N_{2}, N_{3}$ - квадратные матрицы одинакового размера, а 0 - нулевая матрица. Две квадратные матрицы $A_{1}$ и $A_{2}$ наз. к о б о р д а т н ы м и, если матрица



\includegraphics[max width=\textwidth, center]{2023_07_09_78e3c4c8d363b6afcd59g-1(1)}

ная матрица $A$ наз. $\varepsilon$-матрицей, где $\varepsilon=+1$ или-1, если $\operatorname{det}\left(A+\varepsilon A^{\prime}\right)= \pm 1$ Матрица Зейферта всякого $(2 q-1)$-мерного узла является $(-1)^{q}$-матрицей. Для всякого $\varepsilon$ отношение кобордантности является отноше-



\includegraphics[max width=\textwidth, center]{2023_07_09_78e3c4c8d363b6afcd59g-1}

Множество классов эквивалентности обозначается через $G_{\varepsilon}$. Операция прямой суммы определяет в $G_{\varepsilon}$ структуру абелевой групшы. Имеется го м о м о рф и 3 м $Л$ е в и на: $\varphi_{q}: C_{2 q-1} \rightarrow G_{(-1)} q$, к-рый сопоставляет классу У. к. $K$ класс кобордизмов матрицы Зейферта узла $K$. Гомоморфизм Левина $\varphi_{q}$ является изоморфизмом при всех $q \geqslant 3$. Гомоморфизм $\varphi_{2}: C_{3} \rightarrow$ $\rightarrow G_{+}$является мономорфизмом, и его образ является подгруппой индекса 2 в $G_{+}$, состоящей из классов (\begin{CJK}{UTF8}{mj}十\end{CJK}1)-матриц $A$, для к-рых сигнатура матрицы $A+A^{\prime}$ делится на 16. Гомоморфизм $\varphi_{1}: C_{1} \rightarrow G_{-1}$ является әпиморфизмом; его ядро нетривиально.



Для изучения строения груши $G_{+}$и $G_{-}$и построения полной системы инвариантов класса У. к. используется следующая конструкция. И з м е т р и ч е ск о й с т р у к т у р о й над полем $F$ наз. пара $(<,>$; $T$ ), состоящая из невырожденной квадратичной формы $<$, $>$, заданной на конечномерном векторном пространстве $V$ над полем $F$, и ее изометрии $T: V \rightarrow V$. Изометрич. структура $(\langle,>; T)$ наз. к о б о р д а н $\mathrm{T}$ н о й н у л ю, если $V$ содержит вполне изотропное инвариантное относительно $T$ подпространство половинной размерности. Операции ортогональной суммы



$\triangle 16$ Математическая әнц., т. 55





\end{document}